\section{Annexes}

\printbibliography % Output the bibliography

\subsection{Images}
%\subsubsection{Valeur seuil}

\begin{figure}[h]
    \centering
    \begin{minipage}{0.45\textwidth}
        \centering
        \includegraphics[width=\linewidth]{Valeur_seuil_04.png}
		\subcaption{Valeur seuil à $0.4$}
        \label{fig:valeurSeuil1}
    \end{minipage}
    \hfill
    \begin{minipage}{0.45\textwidth}
        \centering
        \includegraphics[width=\linewidth]{Valeur_seuil_061.png}
        \subcaption{Valeur seuil à $0.6$}
        \label{fig:valeurSeuil2}
    \end{minipage}
    \caption{Surface 3D avec une valeur seuil fixée}
    \label{fig:valeurSeuil}
\end{figure}


\begin{figure}[h]
    \centering
    \begin{minipage}{0.45\textwidth}
        \centering
        \includegraphics[width=\linewidth]{mask.png}
		
    \end{minipage}
    \hfill
 
    \caption{Réapplication de Perlin sur une surface 3D}
\end{figure}


%\subsubsection{Fréquence}

\begin{figure}[h]
    \centering
    \begin{minipage}{0.45\textwidth}
        \centering
        \includegraphics[width=\linewidth]{freq_0.png}
		\subcaption{Fréquence $\in \mathbb{N}$}
        \label{fig:freqEntier}
    \end{minipage}
    \hfill
    \begin{minipage}{0.45\textwidth}
        \centering
        \includegraphics[width=\linewidth]{freq_002.png}
        \subcaption{Fréquence à $0.02$}
        \label{fig:freqRegular}
    \end{minipage}
	\caption{Cas classiques de la fréquence}
    \label{fig:freqReg}
\end{figure}

\begin{figure}[h]
    \centering
    \begin{minipage}{0.45\textwidth}
        \centering
        \includegraphics[width=\linewidth]{freq_0001.png}
        \subcaption{Fréquence à $0.001$}
        \label{fig:freqNearZero}
    \end{minipage}
    \hfill
    \begin{minipage}{0.45\textwidth}
        \centering
        \includegraphics[width=\linewidth]{freq_2001.png}
        \subcaption{Fréquence à $2.001$}
        \label{fig:freqAliasing}
    \end{minipage}
    \caption{Cas limites de la fréquence}
    \label{fig:freqLimits}
\end{figure}

%\subsubsection{Octaves}

\begin{figure}[h]
    \centering
    \begin{minipage}{0.45\textwidth}
        \centering
        \includegraphics[width=\linewidth]{octaves1.png}
        \subcaption{Surface avec une seule octave}
        \label{fig:octave1}
    \end{minipage}
    \hfill
    \begin{minipage}{0.45\textwidth}
        \centering
        \includegraphics[width=\linewidth]{octaves3.png}
        \subcaption{Surface avec trois octaves}
        \label{fig:octave3}
    \end{minipage}
    \caption{Comparaison de surfaces avec les octaves}
    \label{fig:octaveReg}
\end{figure}

\begin{figure}[h]
    \centering
    \begin{minipage}{0.45\textwidth}
        \centering
        \includegraphics[width=\linewidth]{octaves5.png}
        \subcaption{Surface avec cinq octaves}
        \label{fig:octave5}
    \end{minipage}
    \hfill
    \begin{minipage}{0.45\textwidth}
        \centering
        \includegraphics[width=\linewidth]{octaves6.png}
        \subcaption{Surface avec six octaves}
        \label{fig:octaves6}
    \end{minipage}
    \caption{Saturation des octaves}
    \label{fig:octaveLimit}
\end{figure}
